\documentclass[11pt]{article}
\usepackage[margin=1in]{geometry}
\usepackage{amsmath,amssymb,amsthm,hyperref}
\newtheorem{theorem}{Theorem}\newtheorem{lemma}[theorem]{Lemma}
\newcommand{\E}{\mathbb E}\newcommand{\R}{\mathbb R}
\newcommand{\relu}{\operatorname{ReLU}}
\newcommand{\one}{\boldsymbol1}
\title{Independent large-width constructions for Gaussian ReLU posteriors}
\author{ChatGPT-6 Astra\thanks{Independent AI research. Sole author. Verification scope and provenance are stated in the accompanying research collection.}}
\date{19 November 2026; independent research draft}
\emergencystretch=3em
\hypersetup{pdftitle={Independent large-width constructions for Gaussian ReLU posteriors},pdfauthor={ChatGPT-6 Astra}}
\begin{document}\maketitle
\begin{abstract}
For a shallow Gaussian ReLU network conditioned on all labels of an
even parity, through hard positive margins or a width-scaled square or
logistic likelihood, we give
an independent sampling construction whose total-variation error tends
to zero as width grows. The constructions use the critical one-neuron
law, a bounded rejection weight on the ordinary neurons, Gaussian
amplitude corrections, and the limiting transverse marks of the
exceptional neurons. They require no unknown evidence, critical mean,
or critical covariance. At each fixed even dimension at least six and
fixed conditioning strength in the condensed phase, their expected number
of elementary proposals and arithmetic operations is linear in width.
The dimension-dependent constant may be very large. This is an
asymptotic real-arithmetic result, with no finite-width error rate or
bit-complexity bound. The proof uses the exact-count and orientation-law
companions; those inputs and this result are not Lean verified.
\end{abstract}

\section{Model and scope}
Fix even $r\ge6$, put $X=\{-1,1\}^r$, $M=2^r$, $K=M/2$, and
$\chi(x)=\prod_i x_i$. Let $\theta_j=(a_j,U_j)$ be independent standard
Gaussian vectors in $\R^{r+2}$ under $P_0^N$, with
\[
 \Phi_x(\theta)=a\relu(U\cdot(1,x)),\qquad
 F_N(x)=N^{-1}\sum_{j=1}^N\Phi_x(\theta_j).
\]
The posterior studied here has density proportional to
\[
 \exp\left[-\frac{N\lambda}{M}\sum_{x\in X}
                                  \ell(\chi(x)F_N(x))\right]
 \quad\text{with respect to }P_0^N,                       \tag{1}
\]
where $\ell(y)=(1-y)^2/2$ or $\ell(y)=\log(1+e^{-y})$.
Every cube label appears in (1); this is not scalar-statistic conditioning.

Write $Z=\E_x\chi(x)\Phi_x$, and set
\[
 \rho=\frac{\binom{r-2}{(r-2)/2}}{2^{r-1}\sqrt r},\quad
 t_*=\rho^{-1},\quad Q_*=M_*^{-1}e^{t_*Z}P_0,\quad
 m_c=\E_{Q_*}Z,\quad \Sigma=\operatorname{Cov}_{Q_*}(\Phi).
\]
The preceding companions prove $M_*<\infty$, $m_c<3/8$, finite
fourth parameter moments, and positive definite $\Sigma$.
Let $s$ minimize $g(y)=t_*y+\lambda\ell(y)$, and suppose $s>m_c$.
Explicitly,
\[
 s=1-t_*/\lambda\quad\text{(square)},\qquad
 s=\log(\lambda/t_*-1)\quad\text{(logistic)}.
\]
For the logistic case the formula requires $\lambda>t_*$; the stronger
assumption $s>m_c$ implies it. Put
\[
                    \gamma=\lambda\ell''(s)/M>0.
\]
The values $s,\gamma,\rho$ are known from the model. Neither $m_c$ nor
$\Sigma$ is needed in the algorithm. The sufficient condition $s>3/8$
can be used when a certified value of $m_c$ is unavailable.

\section{Finite geometry and independent building blocks}
For $z\in X$, define
\[
 v_z=(0,z/\sqrt r),\quad
 \sigma_z=\operatorname{sgn}(G(v_{\one}))\chi(z),\quad
 f_z(x)=\sigma_z\relu(z\cdot x/\sqrt r),
 \qquad G(U)=\E_x\chi(x)\relu(U\cdot(1,x)).
\]
Let $\mathcal B$ contain the orientation sets selecting one representative
of every antipodal pair and obeying
$\sum_{z\in S}\chi(z)z_i=0$ for all $i$. This finite set is nonempty:
a coordinate hemisphere is an example. Order its columns by a fixed
antipodal-pair order and write $A_S=(f_z)_{z\in S}$, an $M$ by $K$
matrix. The parity geometry gives
\[
 \operatorname{rank}A_S=K,\qquad A_S\one=K\rho\chi.
 \tag{2}
\]
In an orthonormal Walsh basis its even-degree block $E$ is invertible
and independent of the choices in $S$. Its remaining nonzero block
consists of singleton characters. Consequently
\[
 A_S^\top A_S\succeq E^\top E,\qquad
 d_E:=\sqrt{\det(E^\top E)}>0,\qquad
 J_S:=\det(A_S^\top A_S)^{-1/2}\le d_E^{-1}.
 \tag{3}
\]
Put $A_S^+=(A_S^\top A_S)^{-1}A_S^\top$ and
$P_S=I-A_SA_S^+$, the orthogonal projection onto the complement of
its column space. In particular $P_S\chi=0$.

Three independent building blocks have finite fixed-dimension expected
proposal cost.
\begin{enumerate}
\item A uniform $S\in\mathcal B$ can be drawn by independent fair choices
of the two orientations in every pair, rejecting unless the displayed
balance equations hold. The acceptance probability is positive. This
method is not claimed to be efficient in $r$.
\item The critical-law rejection construction in the exponential ReLU
condensation companion draws $\theta\sim Q_*$, without knowing $M_*$.
Its acceptance probability is $(r-2)M_*/r\ge(r-2)/r$.
\item At each selected ray, independently draw $v\sim N(0,1)$ and a
transverse mark $t\in v_z^\perp$ with density
\[
 H_r^{-1}e^{-h_z(t)/\rho},\qquad
 h_z(t)=2^{-r-1}\sum_{x:z\cdot x=0}|t\cdot(1,x)|,
 \quad H_r=\int e^{-h_z(t)/\rho}dt.                       \tag{4}
\]
The function $h_z$ is a norm. A product Laplace envelope from norm
equivalence samples (4) at a positive fixed-dimension acceptance rate,
without using $H_r$. Together these marks have probability law
$\mathcal M_z(dv\,dt)$.
\end{enumerate}

\section{The construction}
Take $N>K$. All restarts below redraw the whole proposal.
\begin{enumerate}
\item Draw $S$ uniformly from $\mathcal B$, and independently draw
$n=N-K$ ordinary neurons from $Q_*$. Put
$Y=\sum_{i=1}^n\Phi(\theta_i)$.
\item Accept this pair $(S,\theta_1,\ldots,\theta_n)$ with probability
\[
          d_E J_S\exp\left[-\frac\gamma{2N}\|P_SY\|_2^2\right].
 \tag{5}
\]
The probability lies in $[0,1]$ by (3). Otherwise restart.
\item Draw
$\eta\sim N(0,\gamma^{-1}(A_S^\top A_S)^{-1})$, and set
\[
            c=\frac{Ns}{K\rho}\one-A_S^+Y+\sqrt N\,\eta.
 \tag{6}
\]
Draw $K$ independent marks $(v_j,t_j)\sim\mathcal M_{z_j}$.
If any $c_j\le0$, restart. Otherwise put
\[
 A_j=\frac{\sqrt{v_j^2+4c_j}+v_j}{2},\qquad
 R_j=\frac{\sqrt{v_j^2+4c_j}-v_j}{2},\qquad
 a_j^{\rm big}=\sigma_{z_j}A_j,\qquad
 U_j^{\rm big}=R_jv_{z_j}+t_j/\sqrt{c_j}.                \tag{7}
\]
\item Reject if a constructed neuron falls outside its fixed small
critical-ray cone, or if its norm is no larger than the largest
ordinary-neuron norm. Otherwise inject these ordered $K$ exceptional
neurons into a uniformly chosen ordered set of distinct labels among
$1,\ldots,N$, filling the other labels with the ordinary neurons in order.
\end{enumerate}
The critical cones may be any fixed sufficiently small disjoint cones
on which the companion's exact corner coordinates apply. Their choice
does not depend on $N$. Neither (5), (6), nor (7) uses an estimated
critical mean, covariance, or normalizing constant.

\begin{theorem}[Independent approximation in total variation]
Let $\widehat\Pi_N$ be the law output by this construction and $\Pi_N$
the posterior (1). At fixed even $r\ge6$ and fixed $\lambda$ with
$s>m_c$,
\[
                         \|\widehat\Pi_N-\Pi_N\|_{\rm TV}\to0.
 \tag{8}
\]
The first rejection step has a strictly positive limiting acceptance
probability. Conditional on passing it, the subsequent restart
probability tends to zero. Hence the expected number of elementary
proposals and arithmetic operations per independent output is
$O_{r,\lambda}(N)$ as $N\to\infty$, in idealized real arithmetic.
No finite-$N$ error rate, dimension-uniform bound, or bit-complexity
bound is asserted. The construction is not an exact finite-width sampler.
\end{theorem}

\section{Acceptance and the Gaussian identity}
Under $Q_*^n$ the output mean is $m_c\chi$. Its finite second moment
and the central limit theorem give
\[
       N^{-1/2}P_SY\Longrightarrow P_SZ_0,
                           \qquad Z_0\sim N(0,\Sigma).
\]
The integrand in (5) is bounded and continuous. Averaging over the
finite nonempty set $\mathcal B$ shows that its acceptance probability
$p_N$ converges to
\[
 p_\infty=\frac{d_E}{|\mathcal B|}\sum_{S\in\mathcal B}
                     J_S\E e^{-\gamma\|P_SZ_0\|^2/2}>0. \tag{9}
\]
This bounded change of density also preserves all ordinary-law events
of probability tending to one. In particular, after acceptance,
\[
 Y/N\to m_c\chi,\quad
 \max_{i\le n}\|\theta_i\|/\sqrt N\to0,\quad
 c_j/N\to\alpha:=\frac{s-m_c}{K\rho}>0.                 \tag{10}
\]
The second claim follows from the finite critical parameter second
moment and a union bound. The Gaussian $\eta$ and the finitely many
marks are tight. Formula (7) consequently puts each constructed neuron
inside its prescribed cone, with squared norm divided by $N$ tending
to $2\alpha$, and strictly above every ordinary-neuron norm. This
proves the claims about restarts.

There is an exact reason for the determinant in (5). For fixed $S,Y$,
\[
 \|Y+A_Sc-Ns\chi\|^2
   =\|P_SY\|^2+
      \left\|A_S\left(c-\frac{Ns}{K\rho}\one+A_S^+Y\right)\right\|^2.
 \tag{11}
\]
The normalized Gaussian density of (6), multiplied by (5), cancels
the factor $J_S^{-1}=\sqrt{\det(A_S^\top A_S)}$.
Thus, before the negligible final restarts, the joint law of $S$, the
ordinary neurons, and $c\in\R^K$ has density proportional to
\[
 Q_*^n(d\theta_1\cdots d\theta_n)\,
 \exp\left[-\frac\gamma{2N}\|Y+A_Sc-Ns\chi\|^2\right]dc
 \tag{12}
\]
with counting measure on $S\in\mathcal B$. Its normalizing constant is
exactly
\[
 T_N=N^{K/2}(2\pi/\gamma)^{K/2}
                 \frac{|\mathcal B|}{d_E}p_N.           \tag{13}
\]
The independent mark density is then $\prod_j\mathcal M_{z_j}$.

\section{Comparison with the full parameter posterior}
This step uses the precise evidence constant of the orientation-law
companion, as well as its exact exceptional-neuron count. We state the
comparison explicitly because weak convergence of a few observables
would not establish (8), whose ambient dimension grows with $N$.

Changing the prior to $Q_*^N$ leaves the residual likelihood
\[
 W_N(V)=\exp\left[-\frac NM\sum_x
             \{g(\chi(x)V_x/N)-g(s)\}\right],
             \qquad V=\sum_i\Phi(\theta_i).              \tag{14}
\]
Let $Z_N=\E_{Q_*^N}W_N$. The leading exponential part of the original
evidence has already been factored out in this definition.
In the coordinates (7), the critical one-neuron density is, uniformly
for bounded marks and $c/N\to\alpha$,
\[
 q_z(c,v,t)\,dc\,dv\,dt
 =(1+o(1))L_r(N\alpha)^{-\nu}\,dc\,\mathcal M_z(dv\,dt),
 \quad \nu=(r+1)/2,\quad
 L_r=\frac{(2\pi)^{-(r+1)/2}}{2M_*}H_r.                 \tag{15}
\]
Its actual output is $cf_z+e_j$ with
$\|e_j\|_\infty\le C_r\|t_j\|$ inside the cone. Formula (15) is the
exact Jacobian and corner calculation from the companions; it does not
replace the neuron by a point mass on its ray.

For a sample of (12), after restricting to the final good event,
$Y+A_Sc-Ns\chi=O_P(\sqrt N)$ and $\sum_je_j=O_P(1)$.
Taylor expansion of $g$ at its minimum $s$, using
$g''(s)=M\gamma$, therefore gives
\[
 \frac{W_N(Y+A_Sc+\sum_je_j)}
 {\exp[-\gamma\|Y+A_Sc-Ns\chi\|^2/(2N)]}\to1
 \quad\text{in probability}.                           \tag{16}
\]
This applies to both losses; no hard-margin approximation is used.

The exact-count theorem puts posterior mass $1-o(1)$ on configurations
having exactly one exceptional neuron in each antipodal pair, balanced
orientations, all exceptional neurons inside their prescribed cones,
and all their norms larger than every ordinary-neuron norm. Put the
ordinary neurons first and order the exceptional ones by their pairs.
Exchangeability makes the original assignment of exceptional indices
uniform among the $(N)_K$ ordered injections, independently of this
canonical configuration.

The orientation-law evidence formula, together with (9)--(13), is
precisely
\[
       Z_N\sim (N)_K L_r^K(N\alpha)^{-K\nu}T_N.         \tag{17}
\]
For clarity, its Gaussian factor can be written either as
$\sum_S J_S\E e^{-\gamma\|P_SZ_0\|^2/2}$ or as
\[
 \det(I+\gamma\Sigma)^{-1/2}
 \sum_S\det\left[A_S^\top(I+\gamma\Sigma)^{-1}A_S\right]^{-1/2}.
\]
The equality follows by integrating a Gaussian first over its component
in $\operatorname{im}A_S$, and agrees with the companion's formula.

On the common good-coordinate space, divide the density of the true
canonical posterior by that of the proposal (12) with its mark law.
The ratio is the product of the normalizing factor in (17), the $K$
ratios in (15), and the likelihood ratio (16), with the final good-event
normalizations. Each factor tends to one in proposal probability.
The true good event and the proposal good event both have probabilities
tending to one. After normalizing on that event, the density ratio
$R_N$ is nonnegative with integral one. Thus
\[
 \E|R_N-1|=2\E(1-R_N)_+\longrightarrow0.
\]
This proves total-variation convergence of the entire canonical
parameter arrays. Applying the same uniform random injection cannot
increase total variation and recovers the labeled networks. Equation
(8) follows.

Finally, the first-stage acceptance stays bounded below by (9), and
final rejection tends to zero by (10). Drawing $N-K$ critical neurons
uses expected $O_r(N)$ proposals. Evaluating their full cube outputs and
the fixed-dimensional matrices also costs $O_r(N)$ arithmetic operations.
Uniform balanced orientations, the $K$ marks, and the Gaussian vector
have finite cost depending only on $r,\lambda$. This proves the stated
cost bound, with no claim that its hidden dimension constant is small.

\section{What the construction says about the posterior}
The ordinary neurons of this full-label posterior are not asserted to
have a joint product-law approximation. The bounded weight in (5)
retains their collective output fluctuation in the directions that the
exceptional amplitudes cannot compensate. Omitting it would generally
change the limiting joint law. By contrast, the scalar-conditioned
posterior in the critical-tilt companion has an asymptotically independent
product bulk; these are different conditioning problems.

Nor does the construction select a canonical hemisphere. Its orientation
law is the determinant-weighted law in the companion, and every balanced
pattern has positive limiting probability. The method uses independent
whole-network proposals, so its mathematical justification does not rely
on a Markov chain crossing between orientation sectors. Practical
performance at growing dimension and a finite-width accuracy guarantee
remain open.
\section{An independent construction for hard observed margins}
The hard-margin posterior also admits an independent construction, but
its ordinary readouts require a conditional Gaussian draw. Replacing the
hard constraint by the quadratic weight of (12) would be incorrect.
Fix $s>m_c$ and let
\[
 \Pi_N^H=P_0^N(\,\cdot\mid \chi(x)F_N(x)>s\text{ for every }x\in X).
 \tag{18}
\]
All labels are observed. The standard margin $s=1/2$ satisfies the
assumption by the analytic critical-mean bound.

For each $S$, choose an $M$ by $K$ matrix $B_S$ with orthonormal
columns spanning the orthogonal complement of $\operatorname{im}A_S$.
Thus $B_SB_S^\top=P_S$. Such a basis can be fixed once per orientation
pattern. Put $D_\chi=\operatorname{diag}(\chi(x))$.
For a transverse mark define the finite output vector
\[
 \psi_z(t)_x=\sigma_z\left[
   \one_{\{z\cdot x>0\}}\,t\cdot(1,x)
       +\one_{\{z\cdot x=0\}}\relu(t\cdot(1,x))\right]. \tag{19}
\]
Here indicator symbols in (19) denote scalar zero-one values. Inside the
critical cone the exceptional output has the exact formula
\[
              \Phi(\theta_j^{\rm big})=c_j f_{z_j}
                           +q_j\psi_{z_j}(t_j),\qquad
                       q_j=A_j/\sqrt{c_j}.              \tag{20}
\]
In particular replacing $q_j$ by one incurs an $o_P(1)$ output error
when $c_j/N\to\alpha>0$ and the marks are tight.

We need a known positive covariance cutoff. Let
\[
 K_0=\E_{U\sim N(0,I_{r+1})}
       [\phi(U)\phi(U)^\top],\qquad
 \phi(U)_x=\relu(U\cdot(1,x)),
 \qquad \epsilon=\frac{r-2}{4r}\lambda_{\min}(K_0)>0.
 \tag{21}
\]
The Gaussian ReLU kernel formula gives every entry of $K_0$ explicitly.
It is positive definite: an invertible finite evaluation matrix of
ReLU features, and positive Gaussian mass near its incoming vectors,
make its quadratic form strictly positive in every nonzero direction.
Thus (21) is defined without any unknown critical-law quantity.

The following procedure is considered for sufficiently large $N$; in
particular $N\ge M$ is necessary for its covariance matrix to have full
rank. The proof establishes positive acceptance for all sufficiently large
widths. Every rejection restarts the whole proposal.
\begin{enumerate}
\item Draw a uniform balanced orientation set $S$, and draw $n=N-K$
independent incoming vectors from the incoming marginal of $Q_*$.
Let $H$ be the $M$ by $n$ matrix with columns $\phi(U_i)$, let
$m_i=G(U_i)/\rho$, and put
\[
 C=B_S^\top HH^\top B_S,\qquad b=B_S^\top Hm.
\]
Reject if $C\not\succeq\epsilon N I_K$.
\item Draw the $K$ independent marks $(v_j,t_j)\sim\mathcal M_{z_j}$,
and $M$ independent positive slack variables $u_x$ of exponential
rate $(M\rho)^{-1}$. Set
\[
 e_\infty=\sum_j\psi_{z_j}(t_j),\qquad
 d=B_S^\top(D_\chi u-e_\infty).
\]
Accept this proposal with probability
\[
 d_EJ_S\frac{(\epsilon N)^{K/2}}{\sqrt{\det C}}
        \exp[-\tfrac12(d-b)^\top C^{-1}(d-b)].           \tag{22}
\]
The covariance cutoff and (3) make this probability at most one.
\item Draw $\xi\sim N(0,I_n)$ independently, and set
\[
 a=m+\xi+H^\top B_SC^{-1}\{d-B_S^\top H(m+\xi)\}.
 \tag{23}
\]
These are the ordinary-neuron readouts. Put $Y=Ha$ and
\[
              c=A_S^+(Ns\chi+D_\chi u-e_\infty-Y).      \tag{24}
\]
Draw no new marks: construct the exceptional neurons from these $c_j$
and the previously drawn $(v_j,t_j)$ using (7).
\item Reject if some $c_j\le0$, a constructed neuron is outside its
critical cone, an exceptional norm is not larger than all ordinary
norms, or the actual combined network fails a hard margin in (18).
Otherwise apply the uniform ordered injection as in the soft construction.
\end{enumerate}
Formula (23) samples the Gaussian law $N(m,I_n)$ conditional on
$B_S^\top Ha=d$. Its covariance may be singular; this is intentional.
Formula (24) supplies the complementary $K$ degrees of freedom, and
the independent $M$-dimensional slack restores the full parameter
law. No point conditioning of the whole network is imposed on its
final output distribution.

\begin{theorem}[Hard-margin independent approximation]
At each fixed even $r\ge6$ and $s>m_c$, the output law
$\widehat\Pi_N^H$ of this procedure satisfies
\[
                   \|\widehat\Pi_N^H-\Pi_N^H\|_{\rm TV}\to0.
 \tag{25}
\]
The combined probability of passing the covariance and rejection tests
has a strictly positive limit. Conditional on passing them, the final
restart probability tends to zero. The expected idealized arithmetic
and elementary-proposal cost is $O_{r,s}(N)$, with the same absence of
finite-width, dimension-uniform, or bit-complexity claims as before.
The algorithm uses neither $m_c$, $\Sigma$, $M_*$, nor a numerical
estimate of the posterior evidence.
\end{theorem}
\begin{proof}
Let $Q_U$ be the incoming marginal of $Q_*$. Its density relative to the
standard incoming Gaussian is $M_*^{-1}e^{G(U)^2/(2\rho^2)}$.
The critical rejection envelope gives $M_*\le r/(r-2)$, so
\[
 \overline K:=\E_{Q_U}\phi(U)\phi(U)^\top
                      \succeq\frac{r-2}{r}K_0\succeq4\epsilon I_M.
 \tag{26}
\]
Its entries have finite means. The law of large numbers therefore gives
$C/N\to C_S:=B_S^\top\overline K B_S\succeq4\epsilon I_K$;
the covariance rejection probability tends to zero, uniformly over the
finite set of $S$.

Put $T=\operatorname{Cov}_{Q_U}(\phi(U)G(U)/\rho)$, which is finite
by the fourth incoming moment. Its mean vector is $m_c\chi$, killed
by $B_S^\top$. The central limit theorem gives
\[
             b/\sqrt N\Longrightarrow W_S,
                         \qquad W_S\sim N(0,B_S^\top T B_S).
\]
The proposed marks and slacks are tight, so $d/\sqrt N\to0$.
The integrand (22), including the cutoff, is bounded by one, and its
limit has strictly positive expectation for each $S$. Thus the combined
acceptance probability converges to a positive number. The resulting
bounded density change preserves all high-probability incoming and
mark events.

The conditional-Gaussian formula (23) follows by removing the projection
of $\xi$ onto the row space of $B_S^\top H$, then adding the unique
minimum-norm correction to the specified value. Before the correction,
$H(m+\xi)$ is the output of independent critical neurons. Its
$B_S$ component is $O_P(\sqrt N)$, whereas $d=O_P(1)$.
On the covariance-good event, the correction to the readout vector has
squared Euclidean norm
\[
 [d-B_S^\top H(m+\xi)]^\top C^{-1}
 [d-B_S^\top H(m+\xi)]=O_P(1).
\]
Also $HH^\top/N\to\overline K$, so its output correction is
$O_P(\sqrt N)$. Consequently $Y/N\to m_c\chi$, the largest
ordinary parameter norm divided by $\sqrt N$ tends to zero, and
(24) gives $c_j/N\to\alpha>0$. The construction is therefore in
its cones and correctly separates exceptional norms with probability
tending to one. Formula (20) gives
\[
 Y+\sum_j\Phi(\theta_j^{\rm big})
    =Ns\chi+D_\chi u+o_P(1).                           \tag{27}
\]
Because the law of the proposed slacks after acceptance has a density
bounded by a fixed multiple of their product exponential density, their
minimum is bounded away from zero with arbitrarily high probability
by first choosing a sufficiently small positive bound. Equation (27)
therefore proves that the final hard-margin rejection probability also
tends to zero.

We now compare densities, retaining their full parameter dimension.
Gaussian disintegration and the coarea Jacobian of $c\mapsto A_Sc$
give the identity, for any nonnegative function $w$,
\begin{align*}
 &\int N(m,I_n)(da)\int_{\R^K}w(Ha+A_Sc+e_\infty)\,dc\\
 &\quad=J_S\int_{\R^M}w(Ns\chi+D_\chi u)
     \varphi_C\bigl(B_S^\top(D_\chi u-e_\infty)-b\bigr)\,du,
                                                               \tag{28}
\end{align*}
where the readout conditional law for a fixed $u$ is exactly (23).
The shift $Ns\chi$ lies in $\operatorname{im}A_S$, and the diagonal
sign matrix has determinant of absolute value one. These facts explain
why neither produces another Jacobian.

Let $W_N^H$ be the hard residual weight after tilting to $Q_*^N$:
\[
 W_N^H(V)=e^{-\sum_x(\chi(x)V_x-Ns)/(M\rho)}
                         \prod_x\one_{\{\chi(x)V_x>Ns\}}.
\]
Equation (28), the exponential slack density, and (22) show that before
its final negligible restarts the proposal has density proportional to
\[
 Q_*^n(d\theta)\prod_j\mathcal M_{z_j}(dv_j\,dt_j)\,
       W_N^H(Y+A_Sc+e_\infty)\,dc                       \tag{29}
\]
with counting measure on $S$, and the covariance cutoff retained.
Denote its normalizing constant by $T_N^H$.
The Gaussian density in (28), multiplied by $N^{K/2}$ and the cutoff,
is bounded by $(2\pi\epsilon)^{-K/2}$. The joint limit just proved
therefore gives, by bounded convergence,
\[
 T_N^H\sim N^{-K/2}(M\rho)^M(2\pi)^{-K/2}
            \sum_{S\in\mathcal B}J_S
                          \det(B_S^\top\Sigma B_S)^{-1/2}.
 \tag{30}
\]
Here total covariance gives $\Sigma=T+\overline K$. Integrating the
limiting density of a centered Gaussian with covariance $C_S$ against
$W_S$ supplies their covariance sum in (30). The finitely many
orientation choices and the independent normalized marks cause no
additional limiting factor.

The elementary block-determinant identity
\[
 J_S\det(B_S^\top\Sigma B_S)^{-1/2}
   =(\det\Sigma)^{-1/2}
                   \det(A_S^\top\Sigma^{-1}A_S)^{-1/2}
 \tag{31}
\]
identifies (30) with the hard orientation-law constant. Thus its actual
residual posterior normalizer obeys
\[
 Z_N^H\sim (N)_K L_r^K(N\alpha)^{-K\nu}T_N^H.          \tag{32}
\]
On proposal-typical coordinates, (15) gives the ratio of the true big
neuron densities to the mark approximation. Equations (20), (27), and
the strictly positive proposed slacks give
\[
 \frac{W_N^H(Y+\sum_j\Phi(\theta_j^{\rm big}))}
      {W_N^H(Y+A_Sc+e_\infty)}\to1
 \quad\text{in probability}.                           \tag{33}
\]
The numerator event already holds after final acceptance, and the
exponential ratio tends to one because the output difference is
$o_P(1)$.

The exact-count theorem again allows canonical ordering of the true
exceptional neurons, losing posterior mass $o(1)$. On the proposal's
good-coordinate space, the density $R_N$ of the corresponding restricted
true law relative to the proposal tends to one in probability, by
(15), (32), and (33). Its integral is at most one; the true law may
also assign mass outside this coordinate subset. Nevertheless,
$\E\min(R_N,1)\to1$ implies $\E R_N\to1$ and
$\E|R_N-1|\to0$. Thus that omitted true mass tends to zero as well.
This avoids assuming in advance that approximate slacks are positive
under the true law. The common random injection recovers the labeled
network and proves (25).

At fixed dimension, forming $H$, $C$, and the Gaussian conditional
correction costs $O_r(N)$ arithmetic operations; matrix inversions
have dimension $K$, independent of $N$. The other proposals have the
finite costs described earlier. A positive limiting acceptance rate
and negligible final rejection give the asserted expected cost.
\end{proof}

The hard construction retains both the odd-output constraint and the
small positive final slacks. It therefore approximates the original
hard-margin posterior rather than a smoothed replacement. As before,
observing the entire cube and fixing the dimension before the width
limit are material assumptions. The result does not supply a polynomial
algorithm in the parity dimension or a posterior sampler for arbitrary
partial data sets.

\paragraph{A less conservative covariance envelope.}
The scalar cutoff in (21) is convenient for the proof but can lead to a
small acceptance constant. An anisotropic alternative uses only the
known prior kernel. Let $D_O$ be the restriction of $K_0$ to the odd
Walsh subspace, fix $0<\kappa<(r-2)/r$, and put
\[
 \Lambda_S=\kappa B_S^\top K_0B_S,\qquad
 h_0=\kappa^{K/2}d_E\sqrt{\det D_O}.
\]
Replace the covariance test by $C\succeq N\Lambda_S$, and (22) by
\[
 h_0 N^{K/2}\,\frac{J_S}{\sqrt{\det C}}
          e^{-(d-b)^\top C^{-1}(d-b)/2}.
\]
This is again bounded by one. To see it, use (31) with $K_0$ in place
of $\Sigma$ and the lower bound
$A_S^\top K_0^{-1}A_S\succeq E^\top D_E^{-1}E$, where $D_E$ is the
prior even Walsh block. Since $K_0$ is diagonal in Walsh coordinates,
\[
 \frac{J_S}{\sqrt{\det(B_S^\top K_0B_S)}}
       \le \frac1{d_E\sqrt{\det D_O}}.
\]
The strict choice of $\kappa$ and (26) make the covariance test pass
with probability tending to one. The new prefactor is independent of
$S$, so it leaves the target density (29) unchanged. The preceding proof
therefore applies verbatim. This improves a proposal envelope; it does
not provide a dimension-uniform cost bound.

\end{document}
